\begin{center}
\section*{9. Looking Back, and Looking Ahead}
\end{center}

\indent \indent In Sections 1--8 we built up the basic games 0, 1, 2, $\star$, $\uparrow$ and Tiny 1 with pawns,
added the Trebuchet so that kings could join, and used sums of games to solve three puzzles. In the rest of this
paper we continue that analysis, but this time we check every position with a computer. We wrote programs that
compute the exact value of any pawn position on a flat or cylindrical board, that play out real chess with
kings, and that build complete endgame tables. With these we went back over Sections 1--8, and then pushed on
to questions that, as far as we could find, nobody has studied on a board without edges.

\indent Because so much of this work is computational, we mark how strongly each claim was checked:
\begin{itemize}
\item \proved: we have a written proof;
\item \checkedtwo: two programs, written separately from the rules alone, give the same answer;
\item \computedone: one program computed it, and nothing else has checked it yet;
\item \conjecture: we believe it, but it is neither proved nor computed.
\end{itemize}

\indent As in Sections 1--8, white is left and black is right, a passed pawn counts as a win, and kings are
frozen in a Trebuchet so that they count as the zero game. Section 15 asks when that last assumption is
actually true.

\subsection*{9.1 What held up, and what did not}

\indent The first thing we did was re-check every pawn position of Sections 1--8. Each diagram was copied out
again without looking at our earlier answer, and the values were computed from scratch. Most of the paper holds
up: the zero games of Section 2, the $\pm1$ and $\pm2$ games of Section 3, $\star$ and Tiny 1 in Section 4, the
Trebuchet in Section 5, and the c-file and f-file games of Puzzle 7.1 are all correct. A few things are not, and
we list them here so that the reader can trust the rest.

\begin{table}[h]
\centering\small
\begin{tabularx}{\textwidth}{@{}l l l X@{}}
\toprule
Where & We said & The computer says & Why \\
\midrule
Position 2.2 & $0$ & $0$ & correct, but the white pawn on h1 could never stand there in chess \\
Position 3.1 & $1$ & $1$ & correct, but only because of en passant across the a/h seam \\
Position 3.2 & $2$ & $2$ & the value is right, the list of options is not \\
Position 4.2 & $\uparrow$ & $0$ & the diagram is one rank off; with the pawns on e2 and e5 it is $\uparrow$ \\
Position 5.1 & $0$ & mutual zugzwang & correct, and only because of the seam; it is a real zugzwang in chess \\
Puzzle 7.1, a/h files & $1$ & $\{\BIG \mid \star\}$ & this part is hot, not a number \\
Puzzle 7.2, pawns & $\pm1$ & an infinitesimal & the pawn game is infinitely small (mean 0), not a switch \\
Puzzle 7.2, kings & $0$ & not a zugzwang & in real chess white to move draws (Section 15) \\
Puzzle 7.3 & $2$ & already decided & the pawns on a5 and f7 are already passed \\
Section 1 & knight h1 to a1, a2 & a3, b2 & a1 and a2 are extra \emph{king} squares, not knight squares \\
Section 8 & the ring rule & see Section 17 & under the rule as stated nobody can make a first move \\
\bottomrule
\end{tabularx}
\caption*{Table 9.1: Re-checking Sections 1--8 \checkedtwo. None of these answers changes under the stricter
passed-pawn rule of Section 10.}
\end{table}

\indent So the method of Sections 1--8, breaking a position into a sum of smaller games and reading off
numbers, infinitesimals and temperatures, is sound. The mistakes were in individual positions, and the biggest
one, Section 8, is fixed in Section 17.

\subsection*{9.2 What is new}

\indent Here is a short guide to the rest of the paper.
\begin{itemize}
\item \textbf{Section 10} makes the rules exact. En passant turns out to break the usual test for a passed pawn,
  so we fix the test.
\item \textbf{Section 11} computes the value of every small pawn position on boards of 3 to 8 files, flat and
  cylindrical, and finds which values only a cylinder can make.
\item \textbf{Section 12} finds Elkies' smallest values, $1/8$, $\uparrow^{2}\star$ and Tiny 1, in positions
  that can really happen in a game.
\item \textbf{Section 13} is about $\star2$. Elkies asked whether $\star2$ can happen with only kings and pawns.
  We show that it does not happen with six pawns or fewer, but that it \emph{does} come back with seven.
\item \textbf{Section 14} checks these values against real chess, using endgame tablebases.
\item \textbf{Section 15} unfreezes the kings.
\item \textbf{Section 16} asks which pieces can still checkmate when there are no corners or edges.
\item \textbf{Section 17} corrects and finishes Dawson's chess on the cylinder.
\item \textbf{Section 18} tests chess engines on these positions.
\item \textbf{Sections 19 and 20} explain how everything was checked and list what we still do not know.
\end{itemize}
